\subsection{Graded ring associated with a regular local ring}

THEOREM 1. --- Let A be a Noetherian local ring. The following conditions are equivalent:

\begin{enumerate}
\def\labelenumi{(\roman{enumi})}
\item
  A is regular.
\item
  The ideal \(\mathfrak{m}_{\mathbf{A}}\) is generated by a secant subset for A (§ 3, no. 2, def. 1).
\item
  The ideal \(\mathfrak{m}_{\mathrm{A}}\) is generated by a completely secant sequence for A (A, X, p.~157, def. 2).
\item
  Let S be the symmetric algebra of the \(\kappa_{\mathbf{A}}\) -vector space \(\mathfrak{m}_{\mathbf{A}} / \mathfrak{m}_{\mathbf{A}}^{2}\) , and let \(\operatorname{gr}(\mathbf{A}) = \bigoplus_{n \geqslant 0} \mathfrak{m}_{\mathbf{A}}^{n} / \mathfrak{m}_{\mathbf{A}}^{n+1}\) be the graded ring associated with A. The canonical homomorphism \(\gamma\) of S onto \(\operatorname{gr}(\mathbf{A})\) is bijective
\item
  There exists an integer \(r \geqslant 0\) such that one has \(H_{\mathbf{A},\mathfrak{m}_{\mathbf{A}}} = (1 - \mathrm{T})^{-r}\) , that is, \(\mathfrak{m}_{\mathbf{A}} = 0\) if \(r = 0\) and \([\mathfrak{m}_{\mathbf{A}}^{n}/\mathfrak{m}_{\mathbf{A}}^{n+1}:\kappa_{\mathbf{A}}] = \binom{n + r - 1}{r - 1}\) for every integer \(n \geqslant 0\) if \(r > 0\) .
\item
  We have \(\mathrm{H}_{\mathrm{A},\mathfrak{m}_{\mathrm{A}}} = (1 - \mathrm{T})^{-d}\) with \(d = \dim (\mathrm{A})\)
\end{enumerate}

If these conditions are satisfied, every system of coordinates of A is a completely secant sequence for A.

\begin{enumerate}
\def\labelenumi{(\roman{enumi})}
\setcounter{enumi}{1}
\item
  \(\Rightarrow\) (i): indeed, every secant subset has at most dim(A) elements ( \(\S 3\) , no. 2, th. 1).
\item
  \(\Rightarrow\) (ii): indeed, every completely secant sequence is secant (§ 3, no. 2, corollary of prop. 3).
\end{enumerate}

\((\mathrm{iv}) \Rightarrow (\mathrm{iii}) : \text{let } (x_1, ..., x_r)\) be a sequence of elements of \(\mathfrak{m}_{\mathbf{A}}\) whose classes modulo \(\mathfrak{m}_{\mathbf{A}}^2\) form a basis \((\xi_1, ..., \xi_r)\) of \(\mathfrak{m}_{\mathbf{A}} / \mathfrak{m}_{\mathbf{A}}^2\) over the field \(\kappa_{\mathbf{A}}\) . If property (iv) is satisfied, \(\operatorname{gr}(\mathbf{A})\) is the polynomial algebra \(\kappa_{\mathbf{A}}[\xi_1, ..., \xi_r]\) and the sequence \((x_1, ..., x_r)\) is completely secant (A, X, p.~160, th. 1). This also proves the last assertion of th. 1.

\begin{enumerate}
\def\labelenumi{(\roman{enumi})}
\tightlist
\item
  \(\Rightarrow\) (iv): set \(r = [\mathfrak{m}_{\mathbf{A}} / \mathfrak{m}_{\mathbf{A}}^{2}:\kappa_{\mathbf{A}}]\) . According to formula (6) of § 4, no. 2, the Poincaré series of the graded vector space S over the field \(\kappa_{\mathbf{A}}\) is equal to
\end{enumerate}

\[
\mathrm{P} _ {\mathrm{S}} = \sum_ {n \geqslant 0} \left[ \mathrm{S} ^ {n} (\mathfrak {m} _ {\mathrm{A}} / \mathfrak {m} _ {\mathrm{A}} ^ {2}): \kappa_ {\mathrm{A}} \right] \mathrm{T} ^ {n} = (1 - \mathrm{T}) ^ {- r}.
\]

Suppose that the canonical homomorphism \(\gamma: S \to \mathrm{gr}(A)\) is not bijective. Since \(\gamma\) is surjective, there exists a homogeneous element \(u\) of \(S\) , of degree \(d > 0\) , annihilated by \(\gamma\) . Then one has

\[
\begin{array}{r l} \mathrm {H_ {A,m_ {A}}} = & \mathrm {P_ {S}} - \symrm {P_ {Ker(\gamma)}} \leqslant \mathrm {P_ {S}} - \mathrm {P_ {uS}} = (1 - \mathrm{T} ^ {d}) / (1 - \mathrm{T}) ^ {r} = \\ & = (1 + \mathrm{T} + \dots + \mathrm{T} ^ {d - 1}) / (1 - \mathrm{T}) ^ {r - 1}. \end{array}
\]

According to th. 3 of § 4, no. 4 and lemma 2 of § 4, no. 1, one therefore has dim(A) \textless{} r, and A is not regular.

Let us finally prove the equivalence of conditions (iv) through (vi). Now, (iv) means that we have \(\mathrm{H}_{\mathrm{A},\mathfrak{m}_{\mathrm{A}}}=(1-\mathrm{T})^{-s}\) with \(s=[\mathfrak{m}_{\mathrm{A}}/\mathfrak{m}_{\mathrm{A}}^{2}:\kappa_{\mathrm{A}}]\). Therefore conditions (iv), (v), (vi) mean that we have \(\mathrm{H}_{\mathrm{A},\mathfrak{m}_{\mathrm{A}}}=(1-\mathrm{T})^{-m}\) , with respectively \(m=[\mathfrak{m}_{\mathrm{A}}/\mathfrak{m}_{\mathrm{A}}^{2}:\kappa_{\mathrm{A}}]\) , \(m\geqslant0\) , \(m=\dim(\mathrm{A})\). But if one has \(\mathrm{H}_{\mathrm{A},\mathfrak{m}_{\mathrm{A}}}=(1-\mathrm{T})^{-m}\) , one has \(\dim(\mathrm{A})=m\) by §4, no. 4, th. 3 and \([\mathfrak{m}_{\mathrm{A}}/\mathfrak{m}_{\mathrm{A}}^{2}:\kappa_{\mathrm{A}}]=m\) (since \((1-\mathrm{T})^{-m}=1+m\mathrm{T}+\cdots\) ). The equivalence of conditions (iv) to (vi) follows immediately.

COROLLARY 1. --- Every regular Noetherian local ring is integrally closed, and in particular an integral domain.

Assume A is regular. Then gr(A) is isomorphic to a polynomial algebra in finitely many indeterminates over a field (th. 1, (iv)). Consequently, gr(A) is a Noetherian integrally closed ring (V, § 1, no. 3, cor. 3 of prop. 13), and hence A is integrally closed (V, § 1, no. 4, prop. 15).

We shall see in a later chapter that every regular Noetherian local ring is factorial.

COROLLARY 2.---Let A and B be Noetherian local rings and σ a local homomorphism from B to A. Suppose A is regular and B is complete. For σ to be bijective, it is necessary and sufficient that it induce bijections of κ \(_{B}\) on κ \(_{A}\) and of m \(_{B}\) /m \(_{B}^{2}\) on m \(_{A}\) /m \(_{A}^{2}\) .

The stated condition is obviously necessary.

Suppose conversely that \(\sigma\) induces isomorphisms of \(\mathrm{gr}_{0}(\mathbf{B})\) onto \(\mathrm{gr}_{0}(\mathbf{A})\) and of \(\mathrm{gr}_{1}(\mathbf{B})\) onto \(\mathrm{gr}_{1}(\mathbf{A})\). Since the ring \(\mathrm{gr}(\mathbf{B})\) is generated by \(\mathrm{gr}_{0}(\mathbf{B}) \cup \mathrm{gr}_{1}(\mathbf{B})\) and that \(\mathrm{gr}(\mathbf{A})\) is the symmetric algebra of the vector space \(\mathrm{gr}_{1}(\mathbf{A})\) over the field \(\mathrm{gr}_{0}(\mathbf{A})\) , the homomorphism \(\mathrm{gr}(\sigma)\) is bijective. Consequently, \(\sigma\) is bijective (III, § 2, no. 8, cor. 3 of th. 1).

COROLLARY 3. --- Let \(k\) be a field and A a local Noetherian \(k\)-algebra, whose residue field is equal to \(k\). For A to be regular, it is necessary and sufficient that its completion \(\hat{A}\) be isomorphic to a \(k\)-algebra of formal series \(k[[X_1, ..., X_n]]\) .

This follows from the equivalence of (i) and (iv) in Thm. 1, and from Prop. 11 of III, § 2, no. 9.

